\section*{HISTORICAL NOTE}
\phantomsection
\addcontentsline{toc}{section}{Historical Note}

(Numbers in brackets refer to the bibliography at the end of this note.)

The notion of an arbitrary function was almost unknown at the beginning of the nineteenth century. A fortiori, the idea of studying sets of functions in general and of endowing them with a topological structure did not appear before Riemann's time (see the Historical Note to Chapter I), and it was only towards the end of the nineteenth century that this idea came into systematic use.

Nevertheless, the idea of convergence of a sequence of real-valued functions had been used, more or less consciously, since the beginnings of the infinitesimal calculus. Of course, by convergence we mean here pointwise convergence; other types of convergence could not have been described until the notions of a convergent series and a continuous function had been precisely defined by Bolzano and Cauchy. The latter at first did not recognize the distinction between pointwise convergence and uniform convergence, and believed that he had proved that the sum of any convergent series of continuous functions is continuous ({[}1{]}, (2), vol.~3, p.~120). The error was pointed out almost immediately by Abel, who at the same time showed that every power series is continuous inside its interval of convergence, by a classical argument which, in this particular case, used essentially the notion of uniform convergence ({[}2{]}, pp.~223-224). It remained only to formulate this notion generally, and this was done independently by Stokes and Seidel in 1847-1848, and by Cauchy himself in 1853 ({[}1{]}, (1), vol.~12, p.~30) (*).

Under the influence of Weierstrass and Riemann, the systematic study of the notion of uniform convergence and related questions was developed in the last third of the nineteenth century by the German school (Hankel, du Bois-Raymond) and above all by the Italians; Dini and Arzelà made precise the necessary conditions for the limit of a sequence of continuous functions to be continuous, while Ascoli introduced the fundamental notion of equicontinuity and proved the theorem which characterizes compact sets of continuous functions {[}3{]} (a theorem popularized later by Montel in his theory of ``normal families'', which are relatively compact sets of analytic functions).

On the other hand, Weierstrass himself discovered \([4b]\) the possibility of uniform approximation of a continuous real-valued function in one or more real variables on a bounded set by polynomials. This result immediately aroused lively interest and led to many ``quantitative'' studies (*).

The modern contribution to these questions has been above all to endow them with the full generality of which they are capable, by considering functions whose domain and range are no longer restricted to R or finite-dimensional spaces, and thus placing them in their natural context with the help of general topological concepts. In particular, the theorem of Weierstrass, which had already shown itself to be a powerful weapon in classical analysis, has been extended in recent years by M. H. Stone to much more general situations; developing an idea introduced by H. Lebesgue (in a proof of Weierstrass's theorem) he has shown clearly the important part played in the theory of approximation of continuous real-valued functions by lattices of functions (approximation by ``lattice polynomials'', cf.~§ 4, Proposition 2 and Theorem 2), and has shown also how the generalized Weierstrass theorem has as immediate consequences a whole series of analogous theorems of approximation, which can thus be grouped together in a much more coherent fashion. We have more or less followed his exposition {[}5{]}.

